\documentclass[UTF8]{ctexart}

%\RequirePackage{amsmath} \RequirePackage{amssymb}
\usepackage{amscd,amsthm,amsfonts,amssymb,amsmath}
\usepackage[colorlinks=true]{hyperref}
%\usepackage{fullpage}
\usepackage[all]{xy}
\usepackage{mathrsfs}
\usepackage{tikz,mathpazo}
%\usetikzlibrary{quotes,angles}
%\usetikzlibrary{calc}
%\usetikzlibrary{decorations.pathreplacing}
\usepackage{bm}
\usepackage{geometry}
\usepackage{xcolor}
\usepackage{eso-pic}
\usepackage{CJK}
\usepackage{arydshln}
\usepackage{mathptmx}
\usepackage[11pt]{moresize}



\newcommand{\BC}{{\mathbb {C}}}

\newcommand{\BN}{{\mathbb {N}}}
\newcommand{\BQ}{{\mathbb {Q}}}
\newcommand{\BR}{{\mathbb {R}}}
\newcommand{\BZ}{{\mathbb {Z}}}
\newcommand{\bv}{{\bm v}}
\newcommand{\bw}{{\bm w}}
\newcommand{\bu}{{\bm u}}
\newcommand{\bx}{{\bm x}}
\newcommand{\by}{{\bm y}}
\newcommand{\rref}{\mathrm{rref}}

\newcommand{\bb}{{\bm b}}
\newcommand{\Coker}{{\mathrm{Coker}}}



\newcommand{\End}{{\mathrm{End}}}
\newcommand{\GL}{{\mathrm{GL}}}

\newcommand{\SO}{{\mathrm{GSO}}}
\newcommand{\Hom}{{\mathrm{Hom}}}

\renewcommand{\Im}{{\mathrm{Im}}}

\newcommand{\Isom}{{\mathrm{Isom}}}
\newcommand{\Id}{{\mathrm{Id}}}



\newcommand{\Ker}{{\mathrm{Ker}}}



\newcommand{\rank}{{\mathrm{rank}}}


\newcommand{\SU}{{\mathrm{SU}}}
\newcommand{\Sym}{{\mathrm{Sym}}}

\newcommand{\sub}{{\mathrm{sub}}}

\newcommand{\tr}{{\mathrm{tr}}}

\newcommand{\matrixx}[4]{\begin{pmatrix}
		#1 & #2 \\ #3 & #4
	\end{pmatrix} }
	
	
	
	\newcommand{\wt}{\widetilde}
	\newcommand{\wh}{\widehat}
	\newcommand{\pp}{\frac{\partial\bar\partial}{\pi i}}
	\newcommand{\pair}[1]{\langle {#1} \rangle}
	\newcommand{\wpair}[1]{\left\{{#1}\right\}}
	\newcommand{\intn}[1]{\left( {#1} \right)}
	\newcommand{\norm}[1]{\|{#1}\|}
	\newcommand{\sfrac}[2]{\left( \frac {#1}{#2}\right)}
	\newcommand{\ds}{\displaystyle}
	\newcommand{\ov}{\overline}
	\newcommand{\incl}{\hookrightarrow}
	\newcommand{\lra}{\longrightarrow}
	\newcommand{\lla}{\longleftarrow}
	\newcommand{\ra}{\rightarrow}
	\newcommand{\imp}{\Longrightarrow}
	\newcommand{\lto}{\longmapsto}
	\newcommand{\bs}{\backslash}
	
	
	
	
	\newtheorem{thm}{定理}
	\newtheorem{cor}[thm]{推论}
	\newtheorem{lem}[thm]{引理}
	\newtheorem{prop}[thm]{命题}
	\newtheorem{defn}[thm]{定义}
	\newtheorem{ques}[thm]{问题}
	\newtheorem{eg}[thm]{例题}
	\newtheorem{ex}[thm]{习题}
	\newtheorem{rmk}[thm]{注释}
		\newtheorem{ntn}[thm]{记号}

	\newcommand{\sk}{\medskip}\newcommand{\bsk}{\bigskip}
	\newcommand{\s}{\sk\noindent}\newcommand{\bigs}{\bsk\noindent}
	
	
	%accented words

	\def\mat(#1,#2,#3,#4){
		\left[\begin{array}{cc}
			#1 & #2 \\ #3 & #4\end{array}
		\right]
	}
		\def\clm(#1,#2){
		\left[\begin{array}{c}
			#1 \\ #2 \end{array}
\right]
	}
		\def\clmn(#1,#2){
		\left[\begin{array}{c}
			#1 \\ \vdots\\ #2 \end{array}
\right]
	}
		\def\clmt(#1,#2,#3){
		\left[\begin{array}{c}
			#1 \\ #2  \\ #3\end{array}
\right]
	}
		\def\clmf(#1,#2,#3){
		\left[\begin{array}{c}
			#1 \\ #2  \\ \vdots \\ #3\end{array}
\right]
	}


		
\begin{document}

\title{习题课材料（四）简要解答}

\date{}

\maketitle

\vspace{-1cm}




\noindent{\Large\textbf{注: 带 $\heartsuit  $号的习题有一定的难度、比较耗时, 请量力为之.}}

\medskip

\noindent{\bfseries 记号}: 如不加说明，我们只考虑实矩阵。对于矩阵$A$, 它的四个基本子空间是列空间$C(A)$, 零空间$N(A)$, 行空间$C(A^T)$和$A^T$的零空间$N(A^T)$。
\medskip

\begin{ex}
		给定$V,W\subset \mathbb{R}^n$。若$\dim V+\dim W>n$，则存在$x\ne0$且$x\in V\cap W$。
\end{ex}

\begin{proof}[证明]
	取$V$的一组基$v_1,\dots,v_p$和$W$的一组基$w_1,\dots,w_q$，考察方程$k_1v_1+\dots+k_pv_p+k_{p+1}w_1+\dots+k_{p+q}w_q=0$。
	由于$\dim \mathbb{R}^n=n$而$p+q>n$，因此上述方程一定有非零解。
	容易验证$u=k_1v_1+\dots+k_pv_p=-k_{p+1}w_1-\dots-k_{p+q}w_q$满足条件。
\end{proof}


\begin{ex}
		$A$是$10$阶方阵，$A^2=0$。则$\rank(A)\le5$。
\end{ex}

\begin{ex}
设$(1,0,0,0,0)^T,(0,1,1,0,0)^T$和$(0,1,1,1,0)^T$构成了$N(A)$的一组基，而$A$为五阶方阵，求$\rref(A)$。 是否可求出$C(A)$, $C(A^T)$和$N(A^T)$的一组基?
\end{ex}
\begin{proof}[提示]
	$C(A^T)$的一组基可求，$C(A),N(A^T)$的不可求。
\end{proof}


\begin{ex}
\begin{enumerate}
\item 设$A$是一个$m\times n$的实矩阵，证明$N(A^TA) = N(A)$。
\item $A$如上。证明$C(A^TA ) = C(A^T)$。
\item 证明$A^TA$可逆当且仅当$A$为列满秩矩阵。
\item 如果$A$是一个$m\times n$的复矩阵，上述命题是否成立？
\end{enumerate}
\end{ex}

\begin{proof}[提示]
1. 若$A^T A x = 0$则$0 = x^T A^T A x = (Ax)^T Ax,$ 由于$Ax$为实向量，这等价于$Ax=0.$

4. 命题对复矩阵不正确。
\end{proof}


\begin{ex}
		设$V$为向量空间，$a_1,\dots,a_n$为$V$中线性无关的向量，证明当且仅当$n$为奇数时，
	$a_1+a_2,a_2+a_3,\dots,a_{n-1}+a_n,a_n+a_1$时线性无关。
\end{ex}



\begin{ex}
	(Steinitz替换定理)
	$a_1,\dots,a_r$线性无关，可用$b_1,\dots,b_s$线性表示，则
	\begin{enumerate}
		\item $r\le s$；
		\item 可以选择$b_1,\dots,b_s$中的$r$个向量换成$a_1,\dots,a_r$，得到的新的向量组与$b_1,\dots,b_s$线性等价。
	\end{enumerate}
\end{ex}


\begin{ex}
	设$A$是$n$阶方阵且$f(x)=a_0+a_1x+\dots+a_nx^n$为$\mathbb{R}$上一多项式，则定义$f(A)=a_0I_n+a_1A+\dots+a_nA^n$。
	已知多项式$f$满足$f(0)=0$，证明对任意方阵$A$，$\rank f(A)\le \rank(A)$。
\end{ex}

\begin{ex}
	证明：$\rank \begin{bmatrix}
		A & b \\
		b^T & 0\\
	\end{bmatrix}=\rank(A)$是$Ax=b$有解的充分不必要条件。
\end{ex}



\begin{ex}[ $\heartsuit  $]
	证明：反对称矩阵的秩是偶数。
\end{ex}
\begin{proof}[提示]
设矩阵为$A=\begin{bmatrix}
		0 & -\alpha_1 & - a_1^T\\
		\alpha_1 & 0 & -a_2^T\\
		a_1 & a_2 & A_3
	\end{bmatrix}$，其中$A_3$为反对称阵。若$\alpha_1=0,a_1=0$，则问题归结为低一阶的反对称矩阵。
	若不然，利用初等行变换将$A$化成$\begin{bmatrix}
		0 & -\alpha_1 & - a_1^T\\
		\beta_1 & * & *\\
		0 & * & *
	\end{bmatrix}$，其中$\beta_1\ne0$，再用同样的初等列变换化成$\begin{bmatrix}
		0 & -\beta_1 & 0\\
		\beta_1 & 0 & - b_2^T\\
		0 & b_2 & B_3
	\end{bmatrix}$。利用初等行变换化成$\begin{bmatrix}
		0 & -\beta_1 & 0\\
		\beta_1 & 0 & -b_2^T\\
		0 & 0 & B_3
	\end{bmatrix}$，再用同样的初等列变换化成$\begin{bmatrix}
		0 & -\beta_1 & 0\\
		\beta_1 & 0 & 0\\
		0 & 0 & B_3
	\end{bmatrix}$，则$\rank(A)=\rank(B_3)+2$，而$B_3$是低两阶的反对称矩阵。

	不断降阶，只需考虑一阶和二阶反对称矩阵，此时结论容易看出。
\end{proof}


\begin{ex}[ $\heartsuit  $]
	设$A$是可逆实反对称矩阵，$b\in \mathbb{R}^n$。证明下列等式成立。
	\begin{enumerate}
		\item $\rank(A+bb^T)=n$。
		\item $\rank\begin{bmatrix}
				A & b\\
				b^T & 0\\
		\end{bmatrix}=n$。
	\end{enumerate}
\end{ex}
\begin{proof}[提示]
	\begin{enumerate}
		\item 考虑方程$(A+bb^T)x=0$。则$x^T(A+bb^T)x=0$。但$x^TAx=0$，因此$(b^Tx)^2=0$，即$b^Tx=0$，于是$Ax=0$。
			而$A$可逆，故$x=0$，因此$A+bb^T$可逆。
			\item 利用初等行列变换，$\begin{bmatrix}
				A & b\\
				b^T & 0\\
		\end{bmatrix}$可以化成$\begin{bmatrix}
		A & 0 \\ 0 & -b^TA^{-1}b
	\end{bmatrix}$。$A$反对称，则$A^{-1}$反对称，得$b^TA^{-1}b=0$。因此原矩阵的秩和$A$相同。
	\end{enumerate}
\end{proof}



\end{document}
